\begin{studiumworked}{M11. Mitad de profundidad frente a mitad de volumen · Intermedio}
\textbf{Enunciado.} Una geometría cilíndrica interior normalizada tiene $R/r_0=5$, longitud fija y extremos inactivos. Determine la razón de área activa a mitad de profundidad disponible y a mitad de volumen consumible.

\textbf{Datos.} Radio exterior relativo: $5$.

\textbf{Incógnitas.} Magnitudes solicitadas y cierre físico.

\textbf{Hipótesis.} Modelo académico con hipótesis del capítulo. Datos originales; no representan hardware ni ensayo.

\textbf{Modelo físico.} Conservación y relaciones constitutivas declaradas.
\[A_b/A_{b0}=r/r_0\]
\[r_{mitad\ profundidad}=(R+r_0)/2\]
\[r_{mitad\ volumen}^2=(R^2+r_0^2)/2\]
\textbf{Desarrollo.}
\begin{enumerate}
\item La profundidad final es $4r_0$; a su mitad el radio es $r_0+2r_0=3r_0$. La razón de área es $3$.
\item Para volumen, se iguala $r^2-r_0^2=(R^2-r_0^2)/2$. Despejar da $r/r_0=\sqrt{13}=3.605551$, también razón de áreas.
\item Usar $R/2=2.5r_0$ no representaría ninguna de las dos mitades pedidas. Tampoco se ha calculado mitad de tiempo porque falta la ley temporal acoplada.
\end{enumerate}
\textbf{Operaciones para contrastar el cálculo.}
\begin{enumerate}

\item Magnitud 1.

\[\fitmath{y_{1}=\frac{5+1}{2}\approx 3\,}\]

\item Magnitud 2.

\[\fitmath{y_{2}=\sqrt{\frac{{5}^{2}+1}{2}}\approx 3.605551\,}\]

\end{enumerate}

\textbf{Resultados.} $A_b/A_{b0}=3$ a mitad de profundidad y $3.605551$ a mitad de volumen.

\textbf{Interpretación.} El resultado corresponde a las hipótesis indicadas; su interpretación requiere conservar fronteras y unidades.

\textbf{Comprobación.} Derivar el volumen consumido respecto a profundidad produce exactamente el área lateral activa.

\textbf{Límites.} Geometría abstracta; no especifica duración, composición ni empuje.
\end{studiumworked}
